    \section{SISTEMAS DE UNIDADES Y CONSTANTES}
    \textbf{Unidades Naturales (QFT):} Se define $\hbar = c = k_B = 1$.
    \\ La dimensión fundamental es la Energía $[E]$ (GeV).
    \\ Análisis dim: $[M]= [\vec{p}] = [\Gamma] = [E]$\;\;\; $[L] = [T] = [E]^{-1}$ \;\;\; $[\sigma] = [E]^{-2}$ \;\;\; $[\tau]= [E]^{-1}$\\
    \textbf{Sistema CGS (H-L):} $\epsilon_0 = \mu_0 = 1\; \; \; \epsilon_0 \cdot \mu_0 = c^{-2}$ \\
    \textbf{U.N. + CGS:} Se denominará Unidades Naturales (U.N.) a la combinación U.N. + CGS. 
    \begin{center}
        \renewcommand{\arraystretch}{1.1}
        \begin{tabular}{|l|c|l|}
            \hline
            \textbf{Concepto}    & \textbf{Símbolo}     & \textbf{Valor / Relación}                 \\ \hline
            Conversión Longitud  & $\hbar c$            & $197.327 \text{ MeV fm}$                  \\ \hline
            Conversión Tiempo    & $1 \text{ s}$        & $1.52 \times 10^{24} \text{ GeV}^{-1} = 3\times 10^8 \text{ m}$    \\ \hline
            Conversión Tiempo $^{-1}$    & $1 \text{ s}^{-1}$     & $6.6 \times 10^{-22} \text{ MeV}$ \\ \hline
            Conversión Masa      & $1 \text{ kg}$       & $5.61 \times 10^{26} \text{ GeV}$         \\ \hline
            Conversión Carga     & $1 \text{ q}$        & $1.9 \times 10^{-18} \text{ C}$         \\ \hline
            Fermi                & $1 \text{ Fm} = 1 \text{ fm}$       & $5.068 \text{ GeV}^{-1}$                  \\ \hline
            Sección Eficaz       & $1 \text{ b (barn)}$ & $10^2  \text{ Fm}^{2} $                        \\ \hline
            Sección Eficaz       & $1 \text{ GeV}^{-2}$ & $0.389 \text{ mb (mili-barns)}$                        \\ \hline
            Cte. Estructura Fina & $\alpha$             & $e^2/4\pi \approx 1/137$                  \\ \hline
            Electrón             & $1 \text{ e}$        & $0.3028$                        \\ \hline
            Ángulo Weinberg      & $\sin^2\theta_W$     & $\approx 0.231$                           \\ \hline
            Constante Fermi      & $G_F$                & $1.166 \times 10^{-5} \text{ GeV}^{-2}$   \\ \hline
            Masa Proton          & $m_p$                & $\approx 1                   \text{ GeV}$ \\ \hline
            Masa Planck          & $M_P$                & $\approx 1.22 \times 10^{19} \text{ GeV}$ \\ \hline
            Magnetón Bohr        & $\mu_B$              & $e/2m_e$                                  \\ \hline
        \end{tabular}
    \end{center}

    \section{MODELO ESTÁNDAR: MATERIA}
    Fermiones de espín $1/2$. Clasificación en 3 generaciones.

    \begin{center}
        \textbf{TABLA I: LEPTONES ($L=1$, Color=0)}
        \begin{tabular}{|c|c|c|c|}
            \hline
            \textbf{Sabor} & \textbf{Masa (aprox)} & \textbf{Q (e)} & \textbf{Vida Media ($\tau$)}     \\ \hline
            $e^-$          & $0.511$ MeV           & $-1$           & Estable                 \\
            $\nu_e$        & $< 0.8$ eV            & $0$            & -                       \\ \hline
            $\mu^-$        & $105.66$ MeV          & $-1$           & $2.2 \times 10^{-6}$ s  \\
            $\nu_\mu$      & $< 0.17$ MeV          & $0$            & -                       \\ \hline
            $\tau^-$       & $1776.86$ MeV         & $-1$           & $2.9 \times 10^{-13}$ s \\
            $\nu_\tau$     & $< 18.2$ MeV          & $0$            & -                       \\ \hline
        \end{tabular}
    \end{center}

    \begin{center}
        \textbf{TABLA II: QUARKS ($B=1/3$, Color=3)}
        \begin{tabular}{|c|c|c|c|}
            \hline
            \textbf{Sabor} & \textbf{Masa ($\overline{MS}$)} & \textbf{Q (e)} & \textbf{Tipo}       \\ \hline
            $u$ (Up)       & $2.16$ MeV                      & $+2/3$         & Isospín $\uparrow$   \\
            $d$ (Down)     & $4.67$ MeV                      & $-1/3$         & Isospín $\downarrow$ \\ \hline
            $c$ (Charm)    & $1.27$ GeV                      & $+2/3$         & Charm $C=1$          \\
            $s$ (Strange)  & $93.5$ MeV                      & $-1/3$         & Strange $S=-1$       \\ \hline
            $t$ (Top)      & $172.5$ GeV                     & $+2/3$         & Top $T=1$            \\
            $b$ (Bottom)   & $4.18$ GeV                      & $-1/3$         & Bottom $B'=-1$       \\ \hline
        \end{tabular}
    \end{center}
    \textbf{Top Quark:} Único quark que decae antes de hadronizar ($\Gamma_t \approx 1.42$ GeV).

    \section{BOSONES DE GAUGE E INTERACCIONES}
    Spin 1 (Vectores), Higgs spin 0. Mediadores de simetrías locales $SU(3)_C\!\times\! SU(2)_L\! \times\! U(1)_Y$

    \begin{center}
        \textbf{TABLA III: MEDIADORES DE FUERZA}
        \begin{tabular}{|c|c|c|c|}
            \hline
            \textbf{Bosón}   & \textbf{Masa (GeV)} & \textbf{Fuerza} & \textbf{Acoplo}        \\ \hline
            Gluón ($g$)      & $0$       & Fuerte          & $\alpha_s \sim O(1)$   \\
            Fotón ($\gamma$) & $0$       & EM              & $\alpha \approx 1/137$ \\
            $W^\pm$          & $80.37$    & Débil           & $g_w$                  \\
            $Z^0$            & $91.19$   & Débil           & $g_w, g'$              \\
            Higgs ($H$)      & $125.20$   & -               & $\lambda$ (Yukawa)     \\ \hline
        \end{tabular}
    \end{center}
    \textbf{Running Coupling:} Las constantes de acoplo dependen de la escala de energía $Q^2$.
    \\ $\alpha_s(Q^2)$ disminuye al aumentar $Q$ (Libertad Asintótica). $\alpha_{em}(Q^2)$ aumenta levemente.

    \section{FÍSICA DE NEUTRINOS (EXTENDIDO)}
    \textbf{Oscilaciones:} Evidencia de masa no nula ($m_\nu \neq 0$).
    \\ Estados de sabor ($\nu_e, \nu_\mu, \nu_\tau$) $\neq$ Estados de masa ($\nu_1, \nu_2, \nu_3$).
    \\ \textbf{Matriz PMNS:} $(\nu_e \;\;\; \nu_\mu \;\;\; \nu_\tau)^T = U_{PMNS} \cdot (\nu_1 \;\;\; \nu_2 \;\;\; \nu_3)^T$.\\[2pt]
    $U_{PMNS}$ contiene 3 ángulos de mezcla ($\theta_{12}, \theta_{23}, \theta_{13}$) y una fase de violación CP ($\delta_{CP}$).

    \begin{center}
        \begin{tabular}{|l|c|}
            \hline
            \textbf{Parámetro}             & \textbf{Valor}               \\ \hline
            $\Delta m_{21}^2$              & $7.53 \times 10^{-5} \text{ eV}^2$        \\ \hline
            $\Delta m_{32}^2$            & $2.455 \times 10^{-3} \text{ eV}^2$ (NO) \\ \hline
            $\Delta m_{32}^2$            & $-2.529 \times 10^{-3} \text{ eV}^2$ (IO) \\ \hline
            $\sin^2 \theta_{12}$    & $0.307$                      \\ \hline
            $\sin^2 \theta_{23}$    & $0.546$ (NO)          \\ \hline
            $\sin^2 \theta_{23}$    & $0.553$ (IO)          \\ \hline
            $\sin^2 \theta_{13}$    & $0.0219 $                        \\ \hline
            $\delta_{CP}$ (Fase CP)        & $1.19\pi$         \\ \hline
        \end{tabular}
    \end{center}
    Normal Ordering (NO): $m_1 < m_2 \ll m_3$.\;\; Inverted Ordering (IO): $m_3 \ll m_1 < m_2$.
    % \\ \textbf{Número de familias:} $N_\nu = 2.996 \pm 0.007$ (ancho del Z en LEP).

    \section{PROPIEDADES MATEMÁTICAS}
$(AB)^{\dagger} = B^\dagger A^\dagger\quad (A+B)^\dagger =A^\dagger +B^\dagger\quad A=A^\dagger \Rightarrow \lambda_A \in \mathbb{R}\quad  A=-A^\dagger \Rightarrow \lambda_A \in i\mathbb{R}$\\
$\text{Tr}(A_1 A_2 \dots A_n) = \text{Tr}(A_n A_1 A_2 \dots A_{n-1})\quad\text{Tr}(P A P^{-1}) = \text{Tr}(A)\quad \text{Tr}(\tau A)=\tau \cdot\text{Tr}(A) $\\
$
\begin{array}{ll}

e^{\hat{X}}e^{\hat{Y}} = e^{\hat{X}+\hat{Y} + \frac{1}{2}[\hat{X},\hat{Y}] + \dots} &e^{\hat{A}}\hat{B}e^{-\hat{A}} = \hat{B} + [\hat{A},\hat{B}] + \frac{1}{2!}[\hat{A},[\hat{A},\hat{B}]]\;\;\;\; \vec{\mu} = \frac{q \vec{\sigma}}{2m} \\[2pt]
    \det (e^A) =e^{\mathrm{Tr(A)}}\;\;\Gamma = \tau^{-1}&  [A,B]= AB-BA\qquad \{A,B\}=AB+BA=\{B,A\}
\end{array}
$\\
$\quad[A,B] = -[B,A]\qquad [A,B+C]=[A,B]+[A,C]\qquad[A,BC]= [A,B]C+B[A,C]$\\
$\qquad [A,[B,C]]+[B,[C,A]] + [C,[A,B]] =0\qquad [A,A^n] = 0\qquad [A,B]^\dagger= [B^\dagger,A^\dagger]$\\
$\qquad[a A, B] = [A,a B] = a [A,B]\qquad \text{autoestado: } Av = \lambda v\qquad \sigma^i \sigma^j = \delta^{ij}\mathbb{I} + i\epsilon^{ijk}\sigma^k$\\
\quad \quad \quad If $[A, B]= 0 \Rightarrow [A, B^{-1}] = 0 \quad $QED:\;\; $T^\mu_{\;\;\alpha} (T^{-1})^\alpha_{\;\;\nu} = (T^{-1})^\mu_{\;\;\alpha} T^\alpha_{\;\;\nu}= \delta^\mu_\nu\qquad $